\section{Lean formalization blueprint}
\label{sec:formalization}

\subsection{The exact checked starting point}

The inspected repository snapshot is
\texttt{00bbe871268c453903499fb6afe07519e87d9631}.  The relevant namespace is
\begin{center}
\texttt{LeanProofs.IPRA.Combinatorics.BinomialIdentities}.
\end{center}
It imports the exact binomial development and proves, without
\texttt{sorry}, the following named results:
\begin{itemize}[leftmargin=2em]
\item \texttt{binomSumFormula}, \texttt{binom\_sum}: the row sum
      $\sum_{k\le n}\binom nk=2^n$;
\item \texttt{centralFormula}, \texttt{central\_binomial}: the central
      binomial inequality used by that module;
\item \texttt{catalanRatioFormula}, \texttt{catalan\_ratio}: the adjacent
      binomial identity;
\item \texttt{catalanFormula}, \texttt{catalan}: the universal closure of
      Equation~\eqref{eq:checked-catalan}.
\end{itemize}
The binomial implementation in
\texttt{Combinatorics/Binomial.lean} defines
\(\binom nk=n!\mathbin{\mathrm{quo}}(k!(n\monus k)!)\), proves exactness by a
bounded remainder sum, and derives Pascal's rule internally.  This is the
right foundation for the present project; redefining binomial coefficients in
ambient \texttt{Nat} would weaken the formal result.

A minimal reuse file begins as follows.
\begin{lstlisting}[language={},caption={Checked base import and Catalan term}]
import IPRADerivations.Combinatorics.BinomialIdentities

namespace LeanProofs.IPRA.Combinatorics.Catalan

open LeanProofs.ConciseIPRA LeanProofs.IPRA
open LeanProofs.IPRA.Combinatorics.Binomial
open LeanProofs.IPRA.Combinatorics.BinomialIdentities
open Syn
open scoped Fm
open ConciseIPRA.PTm (app zero succ)
open LeanProofs.IPRA.PTm (num monus)

abbrev cat {m : Nat} (n : PTm m) : PTm m :=
  monus (binom (n + n) n) (binom (n + n) (succ n))

-- Existing checked theorem:
#check LeanProofs.IPRA.Combinatorics.BinomialIdentities.catalan
\end{lstlisting}
The supplied \texttt{lean/CatalanProveItBlueprint.lean} records this import and
unproved target formulas without inserting axioms or placeholders.

\subsection{Proposed module graph}

\begin{longtable}{@{}p{0.27\textwidth}p{0.66\textwidth}@{}}
\toprule
\textbf{New module} & \textbf{Responsibility and dependencies}\\
\midrule
\endfirsthead
\toprule
\textbf{New module} & \textbf{Responsibility and dependencies}\\
\midrule
\endhead
\path{Catalan/Core.lean} & Define \texttt{cat} by binomial difference;
prove $n=0$ and instantiate the checked \texttt{catalan}; derive factorial,
odd-binomial, product, and ratio equations.  Depends only on
\texttt{BinomialIdentities}.\\
\path{Catalan/Dyck.lean} & Bit-word coding, prefix heights, Dyck Boolean,
reflection at the first negative prefix, first-return split, and Segner.
Depends on pairing, bounded search, and \texttt{Catalan/Core}.\\
\path{Catalan/ProofVariants.lean} & Chung--Feller cut-and-swap with its
backward-scan inverse, marked triangulation expansion/collapse, and the
balanced-string convolution-cancellation proof.  Depends on
\texttt{Catalan/Dyck} and \texttt{Catalan/Bijections}.\\
\path{Catalan/Bijections.lean} & Binary and plane trees, parentheses,
lattice paths, triangulations, matchings, restricted sequences, and tableaux.
Each class is a bounded code with a \texttt{Valid} formula and explicit
\texttt{encode}/\texttt{decode} inverse lemmas.\\
\path{Catalan/FiniteSums.lean} & Two-colored Motzkin, weighted convolution,
alternating transform, square sum, Narayana/hypergeometric identity, and
Riordan transform.  Depends on \texttt{Core/SumLemmas} and Vandermonde.\\
\path{Catalan/Permutations.lean} & Bounded permutation codes,
standardization, $231$ decomposition, stack sorting, finite RSK, and all
length-three patterns.\\
\path{Catalan/Digits.lean} & Base-$p$ digits, digit sum, carries, Legendre,
valuation formula, parity, base-five automaton, and prime Catalans.  Reuses
\texttt{AnalyticNumberTheory/Primes}.\\
\path{Catalan/Hankel.lean} & Ballot triangles, two Gram identities,
triangular determinants, rectangular Cauchy--Binet, maximal-minor lemma, and
uniqueness.  Reuses \texttt{Core/LinearAlgebra}.\\
\path{Catalan/FormalSeries.lean} & Finite coefficient arrays, Cauchy
product, substitution with zero constant term, $C=1+xC^2$, Bessel coefficient
identities, and Lagrange inversion at fixed coefficient.\\
\path{Catalan/Generalized.lean} & Cycle/Raney lemma, forests, rational and
Fuss--Catalan numbers, generalized Jonah identity, convolution, and
super-Catalan integrality.\\
\path{Catalan/CoxeterAudit.lean} & Type-$A$ interval-root posets, the weighted
maximal-parabolic recurrence, and the finite $A_2/B_2$ source counterexamples.
The all-types theorem remains a separate dependency-heavy target.\\
\path{Catalan/RealCodes.lean} & Rational interval codes for $\pi$, square
root, logarithm, exponential, Riemann sums, Wallis, Stirling remainder,
integrals, and asymptotic moduli.  This is the only genuinely new analytic
infrastructure.\\
\bottomrule
\end{longtable}

\subsection{Representative object-language formulas}

The ratio recurrence should be stored without division.  With variables
$0=n$, a target formula is
\begin{lstlisting}[language={}]
abbrev ratioFormula : Fm PTm 1 :=
  ((succ (succ (var 0))) * cat (succ (var 0)))
    ≐ ((num 2 * succ (var 0 + var 0)) * cat (var 0))
\end{lstlisting}
Here $\doteq$ (written as \texttt{U+2250} in source) is the object-language equality constructor used by the
imported ProveIt modules; the theorem target is the universal closure of this
formula in \texttt{Prf PRAAx}.

For Segner, define a two-variable term body
\[
 b(k,n)=\Cat_k\Cat_{n\monus k},
\]
then use the existing primitive-recursive bounded sum symbol
\texttt{sumF}:
\begin{lstlisting}[language={}]
abbrev segnerBody : PTm 2 := cat (var 0) * cat (monus (var 1) (var 0))
abbrev segnerRhs (n : PTm m) : PTm m :=
  ap2 (sumF segnerBody) (succ n) n
abbrev segnerFormula : Fm PTm 1 :=
  cat (succ (var 0)) ≐ segnerRhs (var 0)
\end{lstlisting}
The proof can either formalize first return or show that both sides satisfy the
checked ratio recurrence.  The bijective route is preferable because later
modules reuse the split maps.

A uniform finite-cardinality API should use
\begin{lstlisting}[language={}]
structure BoundedClass where
  bound : Nat -> Nat
  valid : Nat -> Nat -> Bool

def count (A : BoundedClass) (n : Nat) : Nat :=
  (List.range (A.bound n)).foldl
    (fun s z => s + if A.valid n z then 1 else 0) 0
\end{lstlisting}
only in the executable layer.  The object-language layer mirrors these
functions as primitive-recursive symbols and proves a generic bijection lemma
once.  Subsequent combinatorial theorems instantiate that lemma rather than
repeating sum rearrangements.

\subsection{Formalizing reflection}

Use a bit word $w<2^{2n}$.  Define
\begin{align*}
 \operatorname{up}(w,j)&=\sum_{i<j}\operatorname{bit}(w,i),\\
 \operatorname{down}(w,j)&=j\monus\operatorname{up}(w,j),\\
 \operatorname{dyck}(n,w)&\Longleftrightarrow
   \operatorname{up}(w,2n)=n\ \wedge\
   \sum_{j\le2n}[\operatorname{down}(w,j)>
                         \operatorname{up}(w,j)]=0.
\end{align*}
The first bad position is
\[
 \operatorname{firstBad}(w)=
 \min\{j\le2n:\operatorname{down}(w,j)>
                    \operatorname{up}(w,j)\},
\]
implemented by bounded fold with a found flag, not by unbounded minimization.
\texttt{reflectPrefix} toggles bits below that position.  Prove:
\begin{enumerate}[label=(\roman*),leftmargin=2em]
\item a bad balanced word maps to a word with $n+1$ up bits;
\item the first positive position of the image is the same cut;
\item applying the inverse prefix reflection recovers the input;
\item the maps preserve the bit-length bound.
\end{enumerate}
The generic bijection theorem then gives Equation~\eqref{eq:cat-difference}.
All four goals are bounded bit-count arithmetic and fit quantifier-free
induction on the prefix length.

\subsection{Hankel formalization without analytic dependencies}

Define \texttt{evenBallot r k} and \texttt{oddBallot r k} by the two binomial
differences.  The path-cut bijection first proves Equations
\eqref{eq:even-gram} and \eqref{eq:odd-gram}.  Instantiate ProveIt's
\texttt{det\_mul}, \texttt{det\_transpose}, and triangular determinant lemmas
for each encoded dimension $n$.  If the existing determinant interface is
awkward over natural numbers, work over the repository's integer-difference
ring and coerce all nonnegative entries.

For the third determinant, the only extra algebraic theorem is a rectangular
Cauchy--Binet identity.  A smaller alternative, useful for a first milestone,
is fraction-free elimination on
$[\Cat_{i+j+2}]$: the ballot matrix gives an explicit sequence of unimodular
row and column operations ending in
$\operatorname{diag}(1,\ldots,1,n+1)$.  Both routes are primitive recursive;
Cauchy--Binet better mirrors the human proof and generalizes to shifted
Catalan Hankel polynomials already studied elsewhere in ProveIt.

\subsection{Number-theoretic formalization}

The exact factorial module already imports the prime infrastructure.  Add a
fixed-base digit fold returning the triple
\((\text{digit sum},\text{carry count},\text{next carry})\).  Prove the digit
identity for addition by induction on the number of processed digits.  The
valuation formula then consists of:
\begin{enumerate}[label=(\arabic*),leftmargin=2em]
\item Legendre's bounded sum over powers $p^j$;
\item the finite digit-sum evaluation of that sum;
\item exact valuation of Equation~\eqref{eq:factorial};
\item the carry/digit identity;
\item specialization to $p=2$ and $p=5$.
\end{enumerate}
No unique factorization theorem stronger than the existing least-prime-divisor
induction is required.

\subsection{Formal series and real-code layers}

For formal series, use finite prefixes:
\begin{lstlisting}[language={}]
def coeff (a : Nat -> Int) (n : Nat) : Int := a n
def cauchy (a b : Nat -> Int) (n : Nat) : Int :=
  sumRange (n+1) (fun k => a k * b (n-k))
\end{lstlisting}
Every theorem quantifies over one requested coefficient $n$; there is no type
of completed power series in the object language.  Lagrange inversion can be
proved as a family of polynomial coefficient identities using rooted trees,
which is smaller than a formal derivative library.

For analytic claims, define a rational interval code
\begin{lstlisting}[language={}]
structure RealCode where
  lo hi : Nat -> RatCode
  ordered : forall q, lo q <= hi q
  width : forall q, q * (hi q - lo q) <= 1
  nested : ...
\end{lstlisting}
The ellipsis is not Lean code in the supplied scaffold; the report's
mathematical contract is monotone nesting.  Use Machin intervals for $\pi$,
integer bisection for square root, alternating intervals for logarithm, and
Taylor intervals for exponential.  The target asymptotic theorem is not a
statement about Lean's ambient \texttt{Real}; it is a formula saying that the
computed rational intervals obey Equation~\eqref{eq:leading-modulus}.

\subsection{Milestones and expected proof obligations}

A practical completion order is:
\begin{enumerate}[leftmargin=2em]
\item \textbf{Core}: closed forms, ratio, values, and Segner.
\item \textbf{Dyck/reflection}: one generic finite-bijection theorem and the
      path encodings.
\item \textbf{Finite identities}: all four MathWorld sums, product,
      hypergeometric termination, and $k$-fold convolution.
\item \textbf{Digits}: valuation, parity, final-digit rule, and Catalan primes.
\item \textbf{Hankel}: ballot Gram factorizations and all three determinants.
\item \textbf{Generalized}: rational, Fuss, Jonah, and super-Catalan results.
\item \textbf{Analysis}: boundary sum first (it needs no real library), then
      Bessel coefficient identities, integrals, Wallis, and Stirling.
\end{enumerate}

The finite milestones require only modules already present in the inspected
ProveIt dependency cone plus the new object encodings.  The analytic milestone
is larger, but every theorem statement and modulus has been reduced to
primitive-recursive rational data in this report.

\subsection{Status of the supplied Lean artifacts}

\texttt{CatalanExecutable.lean} contains executable natural-number definitions
and decidable small checks.  \texttt{CatalanProveItBlueprint.lean} records the
exact import, the checked theorem name, and definitions of the first target
terms.  The environment used to build this delivery did not contain a Lean
binary, so these new files are not claimed to have been compiled.  They use no
\texttt{sorry}, no new axioms, and no theorem declaration whose body is
missing.  The already existing ProveIt theorem \texttt{catalan} is separately
kernel checked in the repository.
